%% ============================================================
%% CHAPTER 2 — Background and Related Work   (budget: ~9,000 words)
%% Role: six literatures, each closed with the gap this thesis fills.
%% ============================================================
\chapter{Background and Related Work}
\label{ch:background}

\section{LLM Agents and the ReAct Loop}
\label{sec:bg-react}
% Define the loop: reason -> act -> observe. Tool calling via structured
% API messages; the transcript as the loop's state. Cite ReAct
% \citep{yao2023react}, plan-then-execute variants ReWOO
% \citep{xu2023rewoo}, LLM Compiler \citep{kim2023llmcompiler}.
% Close: all assume host-class execution; none discusses where the loop
% physically runs.

\section{Agent Runtimes: the Claw Family}
\label{sec:bg-runtimes}
% The shrink lineage with a comparison table (OpenClaw, nanobot, PicoClaw,
% ZeroClaw; then the 2026 MCU cluster: ESP-Claw, WireClaw, zclaw, MimiClaw).
% Describe what a runtime provides: loop, tools, memory, channels, autonomy.
% Note the OSS MCU systems have no papers and no measurements — engineering
% exists, science does not. That gap is Ch3+5.

\section{Device--Cloud Collaborative Agents}
\label{sec:bg-edgecloud}
% EcoAgent \citep{yi2025ecoagent} as the closest closed-loop precedent
% (phone-class; on-device 2B VLM verification). DCP \citep{yang2026dcp} as
% the opposite architecture (MCU as validated endpoint). Surveys of
% edge-SLM/cloud-LLM collaboration. Pre-LLM lineage: Embedded-BDI
% \citep{embeddedbdi2026}.
% Close: nobody hosts the loop itself on an MCU; verification without any
% on-device model is unaddressed.

\section{Context and Memory Management for Agents}
\label{sec:bg-context}
% Trajectory/state compression: MEM1 \citep{mem1_2025}, StateAct
% \citep{rozanov2024stateact} (note its JSON-hurts finding), constant-context
% formulations \citep{remember2026,historystate2026}, on-device context
% state objects \citep{samsungctx2025}. Cheap baselines: observation
% masking \citep{lindenbauer2025complexity}. Memory benchmarks:
% MemoryAgentBench \citep{memoryagentbench2026}.
% Close: all measured in tokens; none in bytes/joules/RAM; none addresses
% the re-upload structure of a device-hosted loop (Ch6).

\section{Embedded Scripting and Runtime Extensibility}
\label{sec:bg-scripting}
% Embedded interpreters: Lua/eLua, MicroPython, Berry (Tasmota); dynamic
% loading on MCUs (loadable ELF). ESP-Claw's cap\_lua skill spec as
% engineering precedent (SKILL.md + lua_run_script + async jobs) — describe
% precisely, no evaluation exists. Linux agent skills (OpenClaw, Claude
% Code) and their silent OS dependencies.
% Close: LLM-as-author-and-caller of embedded scripts is unstudied (Ch7).

\section{Speculative Execution for LLM Agents}
\label{sec:bg-speculation}
% Speculative decoding \citep{leviathan2023fast} as the origin. The
% 2025--26 server-side cluster: \citep{guan2025dynamic,spagent2025,
% sherlock2025,bpaste2026,idlespec2026,ghosttool2026}. Their three shared
% assumptions (LLM verifier; latency objective; discardable actions).
% Close: all three invert on physical devices (Ch8).

\section{Evaluating Agent Systems}
\label{sec:bg-benchmarks}
% Model-side benchmarks (ALFWorld \citep{shridhar2020alfworld}, WebShop
% \citep{yao2022webshop}, GAIA, tau-bench, SWE-bench). Harness-side:
% Harness-Bench \citep{harnessbench2026} (fix model, vary harness; the
% 23.8pp finding), RealClawBench \citep{realclawbench2026}, LiveClawBench
% \citep{liveclawbench2026}. Note: no benchmark integrates energy or
% safety in primary scoring; none runs on a device runtime.
% Close: motivates the methodology of Ch4.

\section{Summary of Gaps}
\label{sec:bg-gaps}
% Half a page: a table mapping each gap (loop-on-MCU unmeasured; context
% cost in physical units; skills without OS; speculation with irreversible
% actions; device-capable evaluation) to the chapter that addresses it.
